% Part: first-order-logic
% Chapter: completeness
% Section: introduction

\documentclass[../../../include/open-logic-section]{subfiles}

\begin{document}

\iftag{FOL}
      {\olfileid{fol}{com}{int}}
      {\olfileid{pl}{com}{int}}

\olsection{Inleiding: wat de volledigheidsstelling zegt}

De volledigheidsstelling is een van de belangrijkste resultaten over
logica. Je kunt haar op twee manieren formuleren. We zullen bewijzen
dat beide formuleringen hetzelfde zeggen. De eerste gaat over het
verband tussen logisch gevolg en ons !!{derivation}ssysteem. Stel dat
!!a{sentence}~$!A$ logisch volgt uit enkele !!{sentence}s $\Gamma$.
Dan bestaat er ook !!a{derivation} waarmee we
$\Gamma \Proves !A$ vaststellen. Het !!{derivation}ssysteem is dus zo
sterk als mogelijk: het kan elk werkelijk gevolg afleiden, maar bewijst
niets wat niet werkelijk volgt.

De tweede formulering zegt dat er een model bestaat: elke consistente
verzameling !!{sentence}s is vervulbaar, dus waar te maken.
Consistentie is een begrip uit de bewijstheorie. Het betekent dat ons
!!{derivation}ssysteem bepaalde !!{derivation}s niet kan maken. Maar
waarom zou het ontbreken van zulke !!{derivation}s uit~$\Gamma$
garanderen dat er
\iftag{FOL}{!!a{structure}~$\Struct{M}$}{!!{valuation}{~$\pAssign{v}$}
bestaat waarvoor $\iftag{FOL}{\Sat{M}}{\pSat{v}}{\Gamma}$}? Al voordat
de volledigheidsstelling voor het eerst was bewezen---en zelfs voordat
onze huidige !!{derivation}ssystemen bestonden---dacht de grote Duitse
wiskundige David Hilbert dat een consistente wiskundige theorie ook
garandeert dat de objecten waarover zij spreekt bestaan. In een brief
aan Gottlob Frege schreef hij:
\begin{quote}
  Als de willekeurig gegeven axioma's elkaar met al hun gevolgen niet
  tegenspreken, dan zijn ze waar en bestaan de door de axioma's
  gedefinieerde dingen. Dit is voor mij het criterium voor waarheid en
  bestaan.
\end{quote}
Frege was het daar fel mee oneens. De tweede formulering van de
volledigheidsstelling laat zien waarin Hilbert in elk geval gelijk had.
Als de axioma's consistent zijn, bestaat er \emph{een}
\iftag{FOL}{!!{structure}}{!!{valuation}} die ze allemaal waar maakt.

De volledigheidsstelling, en vooral haar bewijs, is om nog meer redenen
belangrijk. Zij heeft verschillende gevolgen, die we deels apart
bespreken. Neem bijvoorbeeld !!a{derivation} die
$\Gamma \Proves !A$ laat zien. Zo'n !!{derivation} is eindig en kan dus
maar eindig veel !!{sentence}s uit~$\Gamma$ gebruiken. De
volledigheidsstelling zegt daarom het volgende. Als $!A$ een gevolg is
van~$\Gamma$, dan volgt die zin al uit een eindige deelverzameling
van~$\Gamma$. Dit heet \emph{compactheid}. Je kunt hetzelfde zo zeggen:
als elke eindige deelverzameling van $\Gamma$ consistent is, dan is
$\Gamma$ zelf consistent.

Via !!{derivation}s volgt de compactheidsstelling uit de
volledigheidsstelling. Er is ook een directere route: gebruik de
\emph{constructie uit het bewijs} van de volledigheidsstelling. Die
constructie begint met een verzameling !!{sentence}s met een bepaalde
eigenschap, namelijk consistentie. Uit die verzameling bouwt zij
!!a{structure} met de gewenste eigenschap: de structuur vervult de
verzameling. Bijna dezelfde constructie bewijst compactheid direct.
Begin dan niet met een consistente verzameling, maar met een ``eindig
vervulbare'' verzameling !!{sentence}s.
\iftag{FOL}{De constructie levert nog meer resultaten op. Elke
  vervulbare verzameling !!{sentence}s heeft bijvoorbeeld een eindig of
  !!{denumerable} model. Dit heet de stelling van
  L\"owenheim--Skolem. In de logica, en soms ook in de filosofie, bouwt
  men vaak op deze manier !!{structure}s uit verzamelingen
  !!{sentence}s.}{}

\end{document}
